% appendix: auction replica. Blocks moved intact for the 9-page NeurIPS
% main-text boundary (2026-07-19): Methods "Auction replica" subsection,
% Results "Auction-feed sharpness" paragraphs, Discussion "The ceiling and
% the replica path" subsection. Content unchanged except section headers.
\section{Auction replica: machinery, measured sharpness, and the path past
the futures floor}
\label{app:replica}

\subsection{Machinery}
\label{app:replica_machinery}

The replica assembles the first tick from per-stock auction previews
instead of futures:
\begin{equation}
\mu_R = \sum_i w_i \left( \frac{\hat P_i}{C_i} - 1 \right)
\mathbf{1}[\,i \in \mbox{first tick}\,].
\label{eq:replica}
\end{equation}

\begin{center}
\small
\begin{tabular}{ll}
\toprule
Symbol & Meaning \\
\midrule
$w_i$ & stock $i$'s prior-close capitalization weight \\
$\hat P_i$ & predicted first print of stock $i$ \\
$C_i$ & stock $i$'s prior close \\
$\mathbf{1}[\cdot]$ & one when stock $i$ opens in time for the first tick \\
\bottomrule
\end{tabular}
\end{center}

Stale stocks contribute zero, which keeps the whole computation in
divisor-free weighted-gap space. A Monte Carlo layer treats the first tick
as a random variable: certain-live names are pooled analytically,
certain-stale names sit at zero, and uncertain names are simulated under a
common slowness factor (a logit shift, $\beta = 1.0$) with per-stock
preview noise of $0.05\%$ for Nasdaq and $0.30\%$ for NYSE names, both
measured. The simulation outputs mean, sigma, $P(\mathrm{up})$, quantiles,
and the pivotal names, those whose inclusion swings $P(\mathrm{up})$
most, scored as $P(\mathrm{up} \mid \mathrm{in}) - P(\mathrm{up} \mid
\mathrm{out})$. A print filter conditions on crosses observed after
9:30:00, and the seed is fixed. Print timing uses venue priors, Nasdaq
$p_{\mathrm{live}} \approx 0.97$ and NYSE $\approx 0.15$, until real cross
timestamps allow a per-stock fit, with a Bayes update for names still
silent.

Inputs are Nasdaq's Net Order Imbalance Indicator and NYSE order
imbalances \citep{nasdaq2026totalview,nyse2026taqimbalances}. The ITCH
parser streams a raw $\sim$3.5~GB gzipped TotalView-ITCH day file and
finishes a full day in 27 seconds. Development ran on free full-day
samples: seven ITCH days from 2019--2020 \citep{nasdaq2026itchsamples}
and NYSE TAQ imbalance samples including April 1--2, 2026
\citep{nyse2026taqimbalances}. Those dates mostly miss the 2026 quirk
days, so they validate parsers, not the days that matter. The replica
machinery is validated on synthetic data only.

\subsection{Measured feed sharpness}
\label{app:replica_sharpness}

The auction feeds were measured on free sample days. On the one full ITCH
day (December 30, 2019; 3{,}086 Nasdaq symbols), the NOII, Nasdaq's
published indicative opening price, deviates from the official open by a
mean 10.4 basis points at 9:29 and 8.3 at 9:29:50, and by 4.7 and 3.2 when
weighted by the top-100 notional names. At index level that is a replica
error near 1.7 points at 9:29 and 0.75 at 9:29:50, against a futures floor
of roughly 6 points: 3--8$\times$ sharper in the final minute.

The NYSE side
is far noisier: on the TAQ sample day (April 1, 2026; 21 NYSE mega-caps) the
median indicative deviation falls from 100 basis points at 9:00 to 43 at
9:29 and 29 at 9:29:50. That is 6--10$\times$ noisier than Nasdaq NOII,
which validates carrying NYSE names at prior close in the replica. The replica
machinery passes its self-tests (Monte Carlo matches the analytic case; a
synthetic quirk correctly reports JPM swinging $P(\text{up})$ by $-0.16$),
but it has never seen real auction data.

\subsection{The ceiling and the replica path}
\label{app:replica_path}

The 96.3\% ceiling is not a law of nature. It is the limit of
market-price data, and the one priced path past it is the auction
replica: assemble the official first tick from the exchanges' own
pre-open feeds instead of inferring it from futures. Nasdaq's NOII,
``the best predictor of the opening price available to the public''
\citep{nasdaq2026totalview}, runs at one-second cadence from 9:28 on
the post-2021 feed. Measured on sample days, it implies an index-level
replica error near 1.7 points at 9:29 and 0.75 at 9:29:50, against the
$\pm 6$-point futures floor: 3--8$\times$ sharper in the final minute.

The expected payoff was computed before spending anything. It moves 9:29
accuracy from roughly 94\% to 97--98\% and quirk days from 2 of 12
strictly called correctly on futures alone to a majority called. On the
tiniest quirk days the replica prices the abandoned 1--6-cent side at its
true 60--80\%: a several-fold expected-value trade even when the winner
is unknowable. These are projections from measured sample-day sharpness,
not backtest results.

There is no cheaper shortcut. A constituent replica built from pre-market
trades adds nothing over the futures gap, stale prints and regression
weights going to zero, so beating futures requires quotes and auction
data. One design question also stays open. Inverse-variance fusion
(Equation~\ref{eq:fusion}) assumes independent voices, but the futures
and replica voices are nested
and the news and gap voices are correlated; the Monte Carlo scenario
shape is discarded at the Gaussian boundary. Scenario-reweighting fusion
is the proposed fix. We state the problem and leave it unsolved.
